\section{One degree assignment controlling all candidate primes}
\label{sec:global-weight-choice}

\noindent\textbf{Status.}
GW1--GW6 prove the local mean, global finite averaging and exact
weighted-loss choice, including their transport through the family.
The additional \texttt{CoverAll} hypothesis
realizes the selected degrees in a family member. Existence of a
family with that coverage is separate. HW follows the source's active
construction, which selects one pair before forming the final family.

Fix $0<\varepsilon<1/52$ and use the scales in GJ2. All thresholds
below are uniform in the family index. The chain length is $m_{\rm ch}$,
whereas the coefficient representation has rank parameter $k$.

\paragraph{GW1: the actual degree boxes and small mean coefficients.}
The size bounds give
\[
 k\le x^6,\qquad m_{\rm ch}\le x^{1/2+\varepsilon/100},\qquad
 |TT|\le2x^{2+8\varepsilon+\varepsilon/100},\qquad
 N_1\le C_e\log x.
\]
Since $r_G$, $K.\texttt{kII}$ and $C_m=302114$ are fixed, eventually
\[
 K.\texttt{kII}\,r_G^2 kx^{-26}\le x^{-12},\qquad
 C_mN_1/x^{13}\le x^{-12}.                         \tag{GW.1}
\]
For example, $K.\texttt{kII}\,r_G^2\le x^8$ and
$C_m C_e\log x\le x$ suffice. The two integer windows are
\begin{align*}
 I_a&=[\lceil a\rceil,\lfloor a+a/(kx^{13})\rfloor]\cap\mathbb Z,\\
 I_b&=[\lceil b'\rceil,
       \lfloor b'+N_*x^\varepsilon\Lambda b'/k\rfloor]\cap\mathbb Z,
 \qquad N_*=2n\delta_1.
\end{align*}
They are nonempty eventually. Indeed $a\ge x^{19}$,
$b'\ge x^6$, $N_*\ge1$ and $\Lambda\ge1$ eventually, so both real
window widths are at least one. The elementary implication
$v-u\ge1\Rightarrow\lceil u\rceil\le\lfloor v\rfloor$ handles rounding.

\paragraph{GW2: the mean loss for each good orbit.}
Let $p$ be a candidate prime and $v$ a good orbit with $N_v$ selected
positions. For $p<p_0$ its deficit is independent of the degrees:
\[
 \kappa_v=\min(1,K.\texttt{kII}\,N_vr_G^2kx^{-26})
 \le N_vx^{-12}.
\]
For $p\ge p_0$, GJ1 and GI8 give nonnegativity and
\[
 \mathbb E_{\rm local}\kappa_v
 \le C_mN_vN_1/x^{13}\le N_vx^{-12}.                 \tag{GW.2}
\]
Here GI5 supplies the exact numerical and band hypotheses; the
average is over the literal two degree boxes on that orbit. Thus
both assertions in the source's \texttt{MeanKappaBound} hold, including
the pointwise small-prime bound.

\paragraph{GW3: pass from local averages to one common probability space.}
Choose degrees independently and uniformly from $I_a$ and $I_b$ at
each point of $TT$, and assign zero outside $TT$. Write $\mathcal B$
for this finite nonempty product. For a fixed orbit the injection in
GJ3 identifies its $N_v$ positions with a subset of $TT$. Every local
degree pair has exactly
\[
 |I_a|^{|TT|-N_v}|I_b|^{|TT|-N_v}
\]
global extensions. Consequently its marginal average is precisely
the average in (GW.2). Since the good-orbit fibres are disjoint,
\[
 \Sigma_p=m_{\rm ch}\sum_{v\ \mathrm{good}}\kappa_v\ge0,
 \qquad
 \mathbb E_{\mathcal B}\Sigma_p
 \le m_{\rm ch}|TT|x^{-12}.                         \tag{GW.3}
\]
No independence between different primes is needed.

\paragraph{GW4: the total weighted exceptional loss is small.}
Put $t=x^{\varepsilon/3}$ and $X'=e^t$. A harmonic-sum estimate gives
\[
 \sum_{p\ \mathrm{candidate}}\frac{\log p}{p}
 \le t\sum_{j=1}^{\lceil e^t\rceil-1}\frac1j
 \le t(1+t)\le2t^2.                                \tag{GW.4}
\]
This uses only the inclusion of candidate primes in the integers
below $X'$. For $s\ge0$,
$\mathbf1_{\{s>x^{-\varepsilon}\}}\le x^\varepsilon s$.
Define
\[
 Z=\sum_{p\ \mathrm{candidate}}\frac{\log p}{p}
   \bigl(\mathbf1_{\{\Sigma_p>x^{-\varepsilon}\}}+\Sigma_p\bigr).
\]
Equations (GW.3)--(GW.4) and the size bounds imply
\[
 \mathbb E Z
 \le8x^{-19/2+(1453/150)\varepsilon}
 =8x^{-9-\gamma},\qquad
 \gamma=\tfrac12-\tfrac{1453}{150}\varepsilon
 >\tfrac{2447}{7800}>0.                            \tag{GW.5}
\]
This retains the source's stronger intermediate harmonic bound and
uses a slightly simpler, weaker set of size estimates; a prime number
theorem or independence across primes is unnecessary.

\paragraph{GW5: select the single degree pair.}
For $32x^{-\gamma}\le1$, some pair in $\mathcal B$ satisfies
$Z\le x^{-9}/4$, because otherwise its finite average would be larger.
For that same pair, simultaneously over all candidate primes,
\[
 \boxed{\quad
 \sum_{p\ \mathrm{candidate}}\frac{\log p}{p}
  \bigl(2\mathbf1_{\{\Sigma_p>x^{-\varepsilon}\}}+4\Sigma_p\bigr)
 \le4Z\le x^{-9}.
 \quad}                                                   \tag{GW.6}
\]
This is the literal \texttt{WReq}, hence \texttt{WeightVectorExists};
nonnegativity and the bound on $Z$ also give \texttt{MeanEx}.
It controls the weighted exceptional set, not every prime separately.

\paragraph{GW6: transfer the choice through the actual family.}
At a fixed $x\in\texttt{Fm.Win}$ choose a member $i_0$ and apply GW5.
Transfer its selected degrees to the common integer labels
$\texttt{Fm.TTn}(x)$. The field \texttt{bTop\_eq}, together with
$\texttt{LaG}=\Lambda$ from \texttt{fills}, identifies its degree
windows exactly with the family's windows.

Every member at this $x$ has the same frame, selected integer labels,
prime classes, Frobenius steps, scale, orbit selector and canonical
deficit function. These are the explicit \texttt{FamCore}/\texttt{FamP}
equalities. Relabel by the frame identification. The definitions of
candidates, good representatives, pulled-back degrees and $\Sigma_p$
then agree term by term. Thus every member carrying the chosen degree
pair satisfies (GW.6). This proves \texttt{ChosenW} directly from
GW1--GW5, without importing the larger \texttt{LinkG} package.

If \texttt{Fm.CoverAll} also holds, it supplies a member carrying
any admissible pair, in particular this one. This proves
\texttt{ChosenWS}; \texttt{coverS} can then move its distinguished
position while preserving the degrees. \texttt{CoverAll} is a
separate property, not a field of \texttt{FamP}. It is an optional
route: HW below reconstructs the actual preselected-pair transport
without requiring it. Neither route infers full arithmetic-family
existence from finite averaging alone.
